%% ============================================================
\section{Problem Setting}\label{sec:overview}
%% ============================================================

\subsection{Target-Hidden Exit Sufficiency}
\label{sec:overview:task}

Fix a verification instance \((P,Q)\), where
\(P=(\phi_{\mathrm{pre}},B,S)\) contains the precondition, loop guard, and body,
and \(Q\) is a post-loop assertion or function postcondition.
Let \(\mathcal R(P)\) be the states reachable at the loop head from an initial
state satisfying \(\phi_{\mathrm{pre}}\).
A predicate \(I\) is \emph{inductive} when
\[
  \phi_{\mathrm{pre}}\Rightarrow I
  \qquad\text{and}\qquad
  \{I\wedge B\}\,S\,\{I\}.
\]
Inductiveness gives a sound over-approximation of \(\mathcal R(P)\), but does
not make it useful for a particular verification task.
We call \(I\) \emph{exit-sufficient} for \(Q\) when
\[
  I\wedge\neg B\Rightarrow Q.
\]

\begin{definition}[Target-hidden loop verification]
The verification framework receives \((P,Q)\), but the policy receives only
\(P\).
It must generate an invariant set \(I\) without observing \(Q\), and succeeds
only when \(I\) is both inductive and exit-sufficient for the original \(Q\).
\end{definition}

The restriction is on information flow, not output syntax: both visible and
hidden systems emit ACSL \texttt{loop invariant} clauses.
In the visible setting, however, \(Q\) may itself be inductive and can simply be
copied into the loop, as in \texttt{NLA\_lipus/10.c}.
Target hiding removes this shortcut.
It also makes binary verification unusable as a training reward, since querying
the held-out target during RL would leak the answer.
\sam{} instead rewards target-free behavioral strength and retains the original
\(Q\) as the sole final acceptance criterion.

\subsection{A Shared Candidate-Processing Operator}
\label{sec:overview:system}

For a rollout group \(\mathcal A=\{A_1,\ldots,A_G\}\) and response-cap
hyperparameter \(K\), each policy response is first capped at \(K\) parsed
invariant clauses; let \(\bar A_i\) denote that
ordered prefix and \(X_0=\bigcup_i \bar A_i\) its normalized, deduplicated
clause union.
Each optional repair round receives only target-free WP verdicts and grows the
pool monotonically,
\[
  X_{t+1}=X_t\cup R_t.
\]
Finally, \(\mathsf H_P(X_t)\) denotes the subset that survives syntax checking
and the Frama-C/WP Houdini fixpoint on the masked program \(P\).
The sequence
\[
  \mathcal A
  \;\longrightarrow\;
  X_0
  \;\longrightarrow\;
  X_1,\ldots,X_m
  \;\longrightarrow\;
  \mathsf H_P(X_m)
\]
is the shared rollout--Houdini core in Figure~\ref{fig:overview}.

Training and inference differ only in how they consume this operator.
Training evaluates individual rollouts, unions, ablations, and repaired pools
to assign generation and refinement credit.  Within a rollout it preserves
response order to distinguish clauses that extend negative-trace coverage from
later clauses that add nothing; across rollouts it uses union ablation only as
positive complementary credit.
Inference performs the same union, pool growth, and Houdini pruning without
constructing reward traces, then restores \(Q\) for final verification.
Training uses sampled positives as a cheap pre-filter, but soundness in both
modes comes from the same Frama-C/WP Houdini implementation.
Thus deployment is independent of the reward service while preserving the
candidate and prover semantics optimized during training.
